\documentclass[12pt,a4paper]{article}
\usepackage[T1]{fontenc}
\usepackage[utf8]{inputenc}
\usepackage[shorthands=off,brazilian]{babel}
\usepackage[a4paper,margin=2.5cm]{geometry}
\setlength{\parindent}{0pt}
\setlength{\parskip}{0.8em}
\pagestyle{empty}

\begin{document}

\begin{center}
{\large\bfseries Declaração de uso de ferramentas de IA}\\[0.4em]
Exercício 3: Fichamento de artigos científicos\\
Projeto em Ciência de Dados IV, Prof. Rodrigo Richard Gomes, 2\textsuperscript{o} semestre de 2026\\
Pedro Dias Soares e Eduardo Velloso Landeiro
\end{center}

\bigskip
Este trabalho foi produzido com uso extenso de inteligência artificial generativa: o Claude
(Anthropic), modelo Claude Opus 5.5, usado pela interface Claude Code em 3 e 4 de outubro de 2026.
A ferramenta localizou os artigos em acesso aberto e baixou parte deles, extraiu o texto dos PDFs,
ampliou tabelas e figuras para a auditoria do item~12, redigiu os fichamentos, o quadro
comparativo, a síntese e a seção de seleção, e conferiu citações diretas, páginas e números contra
os artigos com o script \texttt{revisao/checar\_fichamento.py}. Também foi usada para uma revisão
crítica da entrega e para tornar o texto mais direto. A seleção dos artigos vem da busca feita
pelos integrantes no Exercício~1. Os integrantes são responsáveis pela revisão final e por todo o
conteúdo entregue.
\end{document}
